\origpage{204}{ROBERTSON AND SEYMOUR}
\noindent by (c), as required. If $c_2$ borders a cuff arising from a major cell $c_0$ of $\Gamma_0$, then again $v \in V(\delta(c_1) \cap \delta(c_0))$ and so $v = \gamma(n)$ for some $n \in \tilde{c}_1$ from (c), as required.

\hspace*{12pt}(P6) Let $c_1$ be a cell bordering a cuff $\Theta$, arising from a major cell $c_0$ of $\Gamma_0$. Number the nodes and cells bordering $\Theta$ as $n_1$, $c_1$, $n_2$, $c_2$, ..., $n_k$, $c_k$, $(n_1)$ as before. Let $v \in V(\alpha(c_1) \cap \alpha(-c_1))$, and choose $d \in C(\Gamma) - \{c_1\}$ with $v \in V(\alpha(d))$, bordering $\Theta$ if possible. We must show that $v \in \beta(n_1 +) \cup \beta(n_2 +)$. Suppose first that $d$ does not border $\Theta$. Then $\alpha(d) \subseteq \delta(d_0)$ for some cell $d_0 \ne c_0$ of $\Gamma_0$, and $v \in V(\delta(c_0) \cap \delta(d_0)) \subseteq \{\gamma(n) \colon n \in \tilde{c}_0 \cap \tilde{d}_0\}$, by (c). Choose $n \in \tilde{c}_0 \cap \tilde{d}_0$ with $\gamma(n) = v$. Since $n \in \tilde{c}_0$, $n$ borders $\Theta$, and $d$ can be chosen bordering $\Theta$, a contradiction. Thus $d$ borders $\Theta$; $d = c_j$ say where $2 \le j \le k$. We may assume that $v \ne \gamma(n_1)$, $\gamma(n_2)$ since otherwise $v \in \beta(n_1 +) \cup \beta(n_2 +)$ as required. Since $v \in V(\alpha(c_1))$, we deduce that $v \in V(\kappa(n_1))$. Also, since $v \in V(\alpha(c_j))$, either $v \in V(\kappa(n_j))$ or $v \in V(\kappa(n_{j+1}))$, and the first occurs unless $v = \gamma(n_{j+1})$. Since $v = \gamma(n_{j+1})$ implies $j \ne k$, we may assume (replacing $j$ by $j+1$ if necessary) that $v \in V(\kappa(n_j))$ where $2 \le j \le k$. From (e)(ii), $v \in V(\kappa(n_2) \cup \kappa(n_k))$, since either $n_j = n_2$ or $n_j = n_k$ or $n_1$, $n_2$, $n_j$, $n_k$ are distinct and in order. If $v \in V(\kappa(n_2))$ then $v \in \beta(n_2 +)$ (since $v \ne \gamma(n_2)$) and if $v \in V(\kappa(n_k))$ then $v \in \beta(n_1 +)$ (since $v \ne \gamma(n_1)$). Thus in either case $v \in \beta(n_1 +) \cup \beta(n_2 +)$, as required.

\hspace*{12pt}(P7) Number the nodes and cells around $\Theta$ as $n_1$, $c_1$, ..., $n_k$, $c_k$, $(n_1)$ as before. Let $1 \le e < f < g < h \le k$. Let $v \in \beta(n_e +) \cap \beta(n_g +)$. We must show that $v \in \beta(n_f +) \cup \beta(n_h +)$. Now $v \in V(\kappa(n_e) \cap \kappa(n_g))$, and so $v \in V(\kappa(n_f) \cup \kappa(n_h))$ by (e)(ii). We may assume then that $v \in V(\kappa(n_f))$ without loss of generality. We may also assume that $v \notin \beta(n_f +)$, and so $v \notin V(\kappa(n_{f-1}))$. Hence $f - 1 \ne e$, and so $e$, $f-1$, $g$, $h$ are distinct and in order. Thus
\begin{equation*}
v \in V(\kappa(n_e) \cap \kappa(n_g)) \subseteq V(\kappa(n_{\delm{j}\fix{f}-1}) \cup \kappa(n_h))\fixmark{fix:13.1-P7}
\end{equation*}
\fixnotetext{A misprint: no $j$ is defined here. The indices in order are $e, f-1, g, h$, and the next sentences use $\kappa(n_{f-1})$.}%
and so $v \in V(\kappa(n_h))$. We may assume that $v \notin \beta(n_h +)$, and so $v \notin V(\kappa(n_{h-1}))$. Hence $h - 1 \ne g$, and so $e$, $f-1$, $g$, $h-1$ are distinct and in order. But $v \in V(\kappa(n_e) \cap \kappa(n_g))$ and $v \notin V(\kappa(n_{f-1}) \cup \kappa(n_{h-1}))$, a contradiction.

This completes the verification that $\pi$ is a portrayal of $G \setminus Z$. For its warp, since each $\beta(n)$ has cardinality $\le p_0$ the warp of $\pi$ is at most $p_0 + 1 = p$, as claimed. Finally, we verify that $\pi$ is $\mathcal{T} \setminus Z$-central. Let $c \in C(\Gamma)$. Then the separation $(\alpha(c), \alpha(-c))$ of $G \setminus Z$ has order $|\tilde{c}| \le 3$ if $c$ is internal, and has order $\le 2p$ from (P6), if $c$ is a border cell. Thus in either case its order is at most $\max(3, 2p)$ and hence less than $\theta - z$. But from (f), $(\alpha(-c), \alpha(c)) \notin \mathcal{T} \setminus Z$, and so $(\alpha(c), \alpha(-c)) \in \mathcal{T} \setminus Z$. Thus $\pi$ is $\mathcal{T} \setminus Z$-central. This completes the proof. $\square$

Next we would like to convert the portrayal given by (13.1) to a $z$-redundant one where $z$ is large, so that the theorems of this paper apply to it; our objective now is to prove (13.4) below.
